\documentclass[10pt,a4paper]{amsart}
\usepackage[margin=19mm]{geometry}
\usepackage{amsmath,amssymb,mathtools}
\usepackage[scr=rsfs,cal=euler]{mathalfa}
\usepackage{xfrac}
\usepackage[unicode,hidelinks,bookmarks=false]{hyperref}
\newcommand{\ie}{i.e.\ }
\newcommand{\cf}{cf.\ }
\newcommand{\resp}{respectively\ }
\newcommand{\Subsection}{\subsection}
\providecommand{\nobreakdash}{\nobreak\hbox{-}}
\setlength{\emergencystretch}{2em}
\multlinegap0pt
\usepackage{bm}
\usepackage{bbm}
\usepackage{dsfont}
\usepackage{xspace}
\usepackage{cancel}
\usepackage{mathtools}
\usepackage{amsthm}
\makeatletter
\@ifundefined{theorem}{\newtheorem{theorem}{Theorem}}{}
\@ifundefined{lemma}{\newtheorem{lemma}[theorem]{Lemma}}{}
\@ifundefined{proposition}{\newtheorem{proposition}[theorem]{Proposition}}{}
\makeatother
\title[The Boundedness of the Bilinear Fractional Integrals along C]{The Boundedness of the Bilinear Fractional Integrals along Curves}
\author{Junfeng Li, Haixia Yu, Minqun Zhao}
\date{Source: arXiv:2411.14830v1 (CC BY 4.0)}
\begin{document}
\maketitle

\noindent\emph{Source and licence.} The Boundedness of the Bilinear Fractional Integrals along Curves. Version arXiv:2411.14830v1 (CC BY 4.0), distributed under the Creative Commons Attribution 4.0 International licence, as recorded in the arXiv licence metadata for this exact version and shown on the abstract page https://arxiv.org/abs/2411.14830v1. Full rights notice: licenses/arxiv-2411.14830-CC-BY-4.0.txt.\par\smallskip

\noindent\textit{Editorial scope.} Counted unit: the Proposition whose source label is \texttt{proposition 4.1} in the article (source0.tex lines 812--816, source label \texttt{proposition 4.1}), delivered together with the article's own complete original proof of it (source0.tex lines 818--985, one \texttt{proof} environment of 168 lines). No mathematics is written, repaired or re-derived: statement and proof are the article's own text line for line. Required context retained uncounted: eq:1.1 (lines 322--324); maintheorem (lines 333--343); lemma 2.1 (lines 455--466); eq:3.1 (lines 627--629); proposition 3.1 (lines 637--640); eq:3.4 (lines 649--651); eq:3.4a (lines 658--662); eq:3.5 (lines 672--674); eq:3.50 (lines 681--683); proposition 3.2 (lines 689--692); eq:4.1 (lines 807--809). Every definition, display and label of those lines is retained unchanged. Omitted from this package and not redistributed: 1--321, 325--332, 344--454, 467--626, 630--636, 641--648, 652--657, 663--671, 675--680, 684--688, 693--806, 810--811, 817--817, 986--1545. Nothing inside a retained excerpt is omitted or altered: every retained body is delivered exactly as the article's lines are written, so no omission receipt of any kind accompanies this package. Citations: the retained excerpts carry no citation command at all, so no bibliography entry of the article is transcribed and no reference is redistributed. Reference adaptation: none was needed and none was made; every reference inside the delivered unit resolves inside it, so no unresolved reference or citation is produced. Printed slips kept literal: no printed word, symbol, hypothesis or number of the counted statement or of its proof was altered. Layout and class: the article's own class is \texttt{article} with packages including amssymb, amsmath, amsthm, color, mathrsfs, indentfirst, graphicx, txfonts; the wrapper is the lane's pinned \texttt{amsart} editorial wrapper at 10pt on A4, which restates the theorem-like environments the retained text needs and adds the article's own macro definitions that the retained text needs (none), with extra wrapper packages (bm, bbm, dsfont, xspace, cancel, mathtools). The delivered numbering is the unit's own. Assets: no retained span contains a figure environment, a graphics-inclusion command, a PDF-page inclusion, a table or a TikZ picture; no image file is copied into this package, the package declares no asset dependency and no image file is redistributed with it. Licence: version arXiv:2411.14830v1 of this article is distributed under the Creative Commons Attribution 4.0 International licence, as recorded in the arXiv licence metadata for this exact version and shown on the abstract page https://arxiv.org/abs/2411.14830v1. A full rights notice is delivered at licenses/arxiv-2411.14830-CC-BY-4.0.txt. Characters: every retained excerpt is pure ASCII, so the delivered engine, XeLaTeX with its default Latin Modern text font, reports no missing character.
\begin{align}\label{eq:1.1}
 \left|\frac{t\gamma'(t)}{\gamma(t)}\right|\leq C^{(1)}_2~~~\textrm{and}~~~\left|\frac{t^j\gamma^{(j)}(t)}{\gamma(t)}\right|\geq C^{(j)}_{1}~~\textrm{with}~~j=1,2.
\end{align}

\begin{theorem}\label{maintheorem}
Assume that $\gamma$ satisfies {\bf(H.)} and $\frac{\omega_{\infty,2}-\omega_{\infty,1}}{\omega_{\infty,2}-1}<\frac{\alpha}{1-\alpha}<\frac{\omega_{\infty,2}}{\omega_{\infty,2}-\omega_{\infty,1}+1}$. Then, there exists a positive constant $C$, such that
$$\left\|H_{\alpha,\gamma}(f,g)\right\|_{L^{r}(\mathbb{R})}\leq C \|f\|_{L^{p}(\mathbb{R})}\|g\|_{L^{q}(\mathbb{R})}$$
holds for all $f\in L^{p}(\mathbb{R})$ and $g\in L^{q}(\mathbb{R})$, where $(\frac{1}{p},\frac{1}{q},\frac{1}{r})$ lies in the open convex hull of the following points:
\begin{align*}
&A\left(\alpha,0,0\right),&&B\left(1,0,1-\alpha\right),\\
&C_1\left(0,\frac{\alpha}{\omega_{\infty,1}},0\right),&&C_2\left(0,\frac{\alpha}{\omega_{0,2}},0\right),\\
&D_1\left(0,1,1-\frac{\alpha}{\omega_{\infty,1}}\right),&&D_2\left(0,1,1-\frac{\alpha}{\omega_{0,2}}\right),\\
&E\left(1, \frac{\omega_{\infty,2}-\alpha}{2\omega_{\infty,2}-\omega_{\infty,1} }, \frac{2(\omega_{\infty,2}-\alpha)}{2\omega_{\infty,2}-\omega_{\infty,1} }\right),&&F\left(\frac{\omega_{\infty,2}-\alpha}{2\omega_{\infty,2}-\omega_{\infty,1} }, 1, \frac{2(\omega_{\infty,2}-\alpha)}{2\omega_{\infty,2}-\omega_{\infty,1} }\right).
\end{align*}
\end{theorem}

\begin{lemma}\label{lemma 2.1}
Assume that $\gamma\in C^{1}(\mathbb{R}^+)$ is monotonic on $\mathbb{R}^+$, $\lim\limits_{t\rightarrow 0^+}\gamma(t)=0$ and satisfies that
\begin{align}\label{eq:2.1}
there~ exist~ positive ~constants~ C^{(1)}_1~ and~ C^{(1)}_2 ~such~ that~ C^{(1)}_1\leq\left|\frac{t\gamma'(t)}{\gamma(t)}\right|\leq C^{(1)}_2~ on~ \mathbb{R}^+.
\end{align}
Then, for any $\varepsilon>0$, there exist positive constants $t_{1,\varepsilon}\in(0,1)$ and $t_{2,\varepsilon}\in(1,\infty)$, such that
$$\begin{cases}
  |\gamma(t)|\in(t^{\omega_{0,2}+\varepsilon},t^{\omega_{0,1}-\varepsilon})~ &\textrm{for}~ \textrm{all} ~t\in(0,t_{1,\varepsilon}); \\
  |\gamma(t)|\in(t^{\omega_{\infty,1}-\varepsilon},t^{\omega_{\infty,2}+\varepsilon})~ &\textrm{for}~ \textrm{all} ~t\in(t_{2,\varepsilon},\infty),
 \end{cases}$$
where the definitions of $\omega_{0,1}$, $\omega_{0,2}$, $\omega_{\infty,1}$ and $\omega_{\infty,2}$ can be found in {\bf(H.)}.
\end{lemma}

\begin{align}\label{eq:3.1}
H_{\alpha,\gamma,j}(f,g)(x):=\int_{2^{j-1}}^{2^j}f(x-t)g(x-\gamma(t))\,\textrm{d}t.
\end{align}

\begin{proposition}\label{proposition 3.1} (Non-critical case $j\neq j_0, j_0+1, j_0+2$) Assume that $\gamma\in C^2(\mathbb{R}^+)$ is monotonic on $\mathbb{R}^+$, $\lim\limits_{t\rightarrow 0^+}\gamma(t)=\lim\limits_{t\rightarrow 0^+}\gamma'(0)=0$, $\omega_{0,1}, \omega_{\infty,1}> 1$, and satisfy \eqref{eq:1.1}. Then, there exists a positive constant $C$, such that
$$\left\|H_{\alpha,\gamma,j}(f,g)\right\|_{L^{r}(\mathbb{R})}\leq C 2^j\|f\|_{L^{p}(\mathbb{R})}\|g\|_{L^{q}(\mathbb{R})}$$
 for all $f\in L^{p}(\mathbb{R})$ and $g\in L^{q}(\mathbb{R})$, where $p,q,r$ satisfy $\frac{1}{p}+\frac{1}{q}=\frac{1}{r}$ and $p, q \in [1, \infty]$.
\end{proposition}

\begin{align}\label{eq:3.4}
\left\|H_{\alpha,\gamma,j}(f,g)\chi_{E}\right\|_{L^{\frac{1}{2}}(\mathbb{R})}\leq |E|\left\|H_{\alpha,\gamma,j}(f,g)\right\|_{L^{1}(\mathbb{R})}\lesssim |E|\|f\|_{L^{1}(\mathbb{R})}\|g\|_{L^{1}(\mathbb{R})}.
\end{align}

\begin{align}\label{eq:3.4a}
\left\|H_{\alpha,\gamma,j}(f,g)\right\|_{L^{\frac{1}{2}}(\mathbb{R})} &=\left\|\sum\limits_{d\in\mathbb{Z},~|d|\leq 3}\sum\limits_{k\in\mathbb{Z}}H_{\alpha,\gamma,j}(f \chi_{I_{k}},g\chi_{I_{k+d}})\right\|_{L^{\frac{1}{2}}(\mathbb{R})} \\
&\leq \sum\limits_{d\in\mathbb{Z},~|d|\leq 3} \left[\sum\limits_{k\in\mathbb{Z}} \left\|H_{\alpha,\gamma,j}(f \chi_{I_{k}},g\chi_{I_{k+d}})\right\|^{\frac{1}{2}}_{L^{\frac{1}{2}}(\mathbb{R})}    \right]^2\nonumber\\
&\lesssim 2^j \sum\limits_{d\in\mathbb{Z},~|d|\leq 3} \left[\sum\limits_{k\in\mathbb{Z}} \left(\left\|f \chi_{I_{k}}\right\|_{L^{1}(\mathbb{R})}\left\|g\chi_{I_{k+d}}\right\|_{L^{1}(\mathbb{R})} \right)^{\frac{1}{2}}  \right]^2\lesssim 2^j\|f\|_{L^{1}(\mathbb{R})}\|g\|_{L^{1}(\mathbb{R})}.\nonumber
\end{align}

\begin{align}\label{eq:3.5}
\left\|H_{\alpha,\gamma,j}(f,g)\chi_{E}\right\|_{L^{\frac{1}{2}}(\mathbb{R})}\lesssim \frac{|E|}{|\gamma'(2^{j-1})|}\|f\|_{L^{1}(\mathbb{R})}\|g\|_{L^{1}(\mathbb{R})}.
 \end{align}

\begin{align}\label{eq:3.50}
\left\|H_{\alpha,\gamma,j}(f,g)\right\|_{L^{\frac{1}{2}}(\mathbb{R})}\lesssim2^j\|f\|_{L^{1}(\mathbb{R})}\|g\|_{L^{1}(\mathbb{R})}
\end{align}

\begin{proposition}\label{proposition 3.2} (Critical cases $j= j_0, j_0+1, j_0+2$) Let $\gamma $ be the same as in Proposition \ref{proposition 3.1}. Then, there exists a positive constant $C$, such that
$$\left\|H_{\alpha,\gamma,j}(f,g)\right\|_{L^{r}(\mathbb{R})}\leq C \|f\|_{L^{p}(\mathbb{R})}\|g\|_{L^{q}(\mathbb{R})}$$
for all $f\in L^{p}(\mathbb{R})$ and $g\in L^{q}(\mathbb{R})$, where $p,q,r$ satisfy $\frac{1}{p}+\frac{1}{q}=\frac{1}{r}$, $p, q \in [1, \infty]$ and $(p,q)\neq (1,1)$.
\end{proposition}

\begin{align}\label{eq:4.1}
H_{\alpha,\gamma}(f,g)(x)\lesssim \sum_{j \in \mathbb{Z}}2^{(\alpha-1)j}H_{\alpha,\gamma,j}(f,g)(x).
\end{align}

\begin{proposition}\label{proposition 4.1}
Let $\gamma $ be the same as in Theorem \ref{maintheorem}. Then, there exists a positive constant $C$, such that
$$\left\|H_{\alpha,\gamma}(f,g)\right\|_{L^{r(\varepsilon),\infty}(\mathbb{R})}\leq C \|f\|_{L^{1}(\mathbb{R})}\left\|g\right\|_{L^{2r(\varepsilon)}(\mathbb{R})},$$
where $\frac{1}{r( \varepsilon ) }:=\frac{2(\omega_{\infty,2}-\alpha)}{2\omega_{\infty,2}-\omega_{\infty,1}}-\varepsilon $ with $\varepsilon>0$ small enough.
\end{proposition}

\begin{proof}
For $H_{\alpha,\gamma,j}$ defined in \eqref{eq:3.1}, let $E$ be a measurable set with finite measure, a calculation gives
\begin{align*}
 \left \|H_{\alpha,\gamma,j}(f,g)\chi_{E}\right\|_{L^{1}(\mathbb{R})}
  \leq\|g\|_{L^{\infty}(\mathbb{R})}\int_{E}\int_{2^{j-1}}^{2^j}f(x-t)\,\textrm{d}t\,\textrm{d}x
  \leq|E|\|f\|_{L^{1}(\mathbb{R})}\|g\|_{L^{\infty}(\mathbb{R})}
 \end{align*}
for all $j\in \mathbb{Z}$. This, combined with Propositions \ref{proposition 3.1} and \ref{proposition 3.2}, shows that
  \begin{align}\label{eq:4.2}
  \left\|H_{\alpha,\gamma,j}(f,g)\chi_{E}\right\|_{L^{1}(\mathbb{R})}\lesssim \min\left\{2^j,|E|\right\}\|f\|_{L^{1}(\mathbb{R})}\|g\|_{L^{\infty}(\mathbb{R})}
  \end{align}
for all $j\in \mathbb{Z}$. From \eqref{eq:3.4} and \eqref{eq:3.4a}, we can write
  \begin{align}\label{eq:4.3}
  \left\|H_{\alpha,\gamma,j}(f,g)\chi_{E}\right\|_{L^{\frac{1}{2}}(\mathbb{R})}\lesssim \min\left\{2^j,|E|\right\}\|f\|_{L^{1}(\mathbb{R})}\|g\|_{L^{1}(\mathbb{R})}
  \end{align}
  for all $j<j_0$, and by \eqref{eq:3.5} and \eqref{eq:3.50}, it follows that
   \begin{align}\label{eq:4.4}
  \left\|H_{\alpha,\gamma,j}(f,g)\chi_{E}\right\|_{L^{\frac{1}{2}}(\mathbb{R})}\lesssim \min\left\{2^j,\frac{|E|}{|\gamma'(2^{j-1})|}\right\}\|f\|_{L^{1}(\mathbb{R})}\|g\|_{L^{1}(\mathbb{R})}
  \end{align}
 for all $j>j_0+2$.

Now, for any fixed $f\in L^1(\mathbb{R})$, the operator $H_{\alpha,\gamma,j}(f,g)\chi_{E}$ may be viewed as a linear operator about $g$. Interpolation \eqref{eq:4.2} with \eqref{eq:4.3} implies that
 \begin{align}\label{eq:4.5}
  \left\|H_{\alpha,\gamma,j}(f,g)\chi_{E}\right\|_{L^{\frac{1}{1+\theta_1}}(\mathbb{R})}\lesssim \min\left\{2^j,|E|\right\}\|f\|_{L^{1}(\mathbb{R})}\|g\|_{L^{\frac{1}{\theta_1}}(\mathbb{R})}
  \end{align}
for all $j<j_0$ with $\theta_1\in (0,1)$. From Proposition \ref{proposition 3.2}, reiterating this interpolation process, we also have
\begin{align}\label{eq:4.7}
  \left\|H_{\alpha,\gamma,j}(f,g)\chi_{E}\right\|_{L^{\frac{1}{1+\theta_1}}(\mathbb{R})}\lesssim \min\left\{2^{j_0},2^{{j_0}\theta_{1}}|E|^{1-\theta_{1}}\right\}\|f\|_{L^{1}(\mathbb{R})}\|g\|_{L^{\frac{1}{\theta_1}}(\mathbb{R})}
  \end{align}
for all $j=j_0, j_0+1, j_0+2$. Using the interpolation argument on \eqref{eq:4.2} with \eqref{eq:4.4}, we claim that
\begin{align}\label{eq:4.6}
  \left\|H_{\alpha,\gamma,j}(f,g)\chi_{E}\right\|_{L^{\frac{1}{1+\theta_1}}(\mathbb{R})}\lesssim \min\left\{2^j,\frac{2^{j(1-\theta_{1})}|E|^{\theta_1}}{|\gamma'(2^{j-1})|^{\theta_{1}}},2^{j\theta_1} |E|^{1-\theta_1},\frac{|E|}{|\gamma'(2^{j-1})|^{\theta_1}}\right\}\|f\|_{L^{1}(\mathbb{R})}\|g\|_{L^{\frac{1}{\theta_1}}(\mathbb{R})}
  \end{align}
for all $j>j_0+2$.

Now pick $E_{\lambda}:=\{x\in{\mathbb{R}}:\ |H_{\alpha,\gamma}(f,g)(x)|>\lambda\}$ with nonnegative functions $f, g \in L^1(\mathbb{R}) \cap L^{\infty}(\mathbb{R})$. One may use Chebyshev's inequality and \eqref{eq:4.1}, \eqref{eq:4.5}-\eqref{eq:4.6} to deduce that
\begin{align}\label{eq:4.8}
 \lambda^{\frac{1}{1+\theta_{1}}}|E_{\lambda}|&\lesssim \int_{E_{\lambda}}\left(\sum\limits_{j\in\mathbb{Z}}2^{(\alpha-1)j}H_{\alpha,\gamma,j}(f,g)(x)\right)^{\frac{1}{1+\theta_{1}}}\,\textrm{d}x
\lesssim\sum\limits_{j\in\mathbb{Z}}2^{{\frac{\alpha-1}{1+\theta_{1}}}j}\int_{E_{\lambda}}H_{\alpha,\gamma,j}(f,g)(x))^{\frac{1}{1+\theta_{1}}}\,\textrm{d}x\\ \nonumber
&\lesssim \Delta_{1}+\Delta_{2}+\Delta_{3},
  \end{align}
where
$$\begin{cases}
   \Delta_{1}:=\sum\limits_{j<j_{0}}2^{{\frac{\alpha-1}{1+\theta_{1}}j}}\min\left\{2^j,|E_\lambda|\right\}^{\frac{1}{1+\theta_{1}}}
\|f\|_{L^{1}(\mathbb{R})}^{\frac{1}{1+\theta_{1}}}\|g\|_{L^{\frac{1}{\theta_1}}(\mathbb{R})}^{\frac{1}{1+\theta_{1}}}; \\
   \Delta_{2}:=2^{{\frac{\alpha-1}{1+\theta_{1}}}j_{0}}\min\left\{2^{j_0},2^{{j_0}\theta_{1}}|E_\lambda|^{1-\theta_{1}}\right\}^{\frac{1}{1+\theta_{1}}}\|f\|_{L^{1}(\mathbb{R})}^{\frac{1}{1+\theta_{1}}}\|g\|_{L^{\frac{1}{\theta_1}}(\mathbb{R})}^{\frac{1}{1+\theta_{1}}};\\
 \Delta_{3}:=\sum\limits_{j>j_{0}+2}2^{{\frac{\alpha-1}{1+\theta_{1}}}j}\min\left\{2^j,\frac{2^{j(1-\theta_{1})}|E_\lambda|^{\theta_1}}{|\gamma'(2^{j-1})|^{\theta_{1}}},2^{j\theta_1} |E_\lambda|^{1-\theta_1},\frac{|E_\lambda|}{|\gamma'(2^{j-1})|^{\theta_1}}\right\}^{\frac{1}{1+\theta_{1}}}\|f\|_{L^{1}(\mathbb{R})}^{\frac{1}{1+\theta_{1}}}\|g\|_{L^{\frac{1}{\theta_1}}(\mathbb{R})}^{\frac{1}{1+\theta_{1}}}.
 \end{cases}$$

These three terms are treated separately in the following section. We are ready to calculate $\Delta_{1}$, which can be written as
\begin{align*}
\left[\sum\limits_{j<j_{0},~{|E_{\lambda}|\geq2^j}}2^{{\frac{\alpha}{1+\theta_{1}}}j}+
\sum\limits_{j<j_{0},~{|E_{\lambda}|<2^j}}2^{{\frac{\alpha-1}{1+\theta_{1}}}j}|E_\lambda|^{\frac{1}{1+\theta_{1}}}\right]
\|f\|_{L^{1}(\mathbb{R})}^{\frac{1}{1+\theta_{1}}}\|g\|_{L^{\frac{1}{\theta_1}}(\mathbb{R})}^{\frac{1}{1+\theta_{1}}}.
\end{align*}
Notice that
$$\sum\limits_{j<j_{0},~{E_{\lambda}\geq2^j}}2^{{\frac{\alpha}{1+\theta_{1}}}j}\lesssim
\begin{cases}
 2^{{\frac{\alpha}{1+\theta_{1}}}{j_0}} , & |E_\lambda|\geq2^{j_0}; \\
  |E_\lambda|^{\frac{\alpha}{1+\theta_{1}}}, & |E_\lambda|<2^{j_0},
\end{cases}$$
and
$$\sum\limits_{j<j_{0},~{|E_{\lambda}|<2^j}}2^{{\frac{\alpha-1}{1+\theta_{1}}}j}|E_\lambda|^{\frac{1}{1+\theta_{1}}}\lesssim
\begin{cases}
 0, & |E_\lambda|\geq2^{j_0}; \\
  |E_\lambda|^{\frac{\alpha}{1+\theta_{1}}}, & |E_\lambda|<2^{j_0}.
\end{cases}$$
Hence we have the estimate
\begin{align}\label{eq:4.9}
\Delta_{1}\lesssim
\begin{cases}
  \|f\|_{L^{1}(\mathbb{R})}^{\frac{1}{1+\theta_{1}}}\|g\|_{L^{\frac{1}{\theta_1}}(\mathbb{R})}^{\frac{1}{1+\theta_{1}}}, & |E_\lambda|\geq2^{j_0}; \\
  |E_\lambda|^{\frac{\alpha}{1+\theta_{1}}}\|f\|_{L^{1}(\mathbb{R})}^{\frac{1}{1+\theta_{1}}}\|g\|_{L^{\frac{1}{\theta_1}}(\mathbb{R})}^{\frac{1}{1+\theta_{1}}}, & |E_\lambda|<2^{j_0}.
\end{cases}
\end{align}

Now we turn to the second term $\Delta_{2}$, it is easy to see that
\begin{align}\label{eq:4.10}
\Delta_{2}\lesssim
\begin{cases}
   \|f\|_{L^{1}(\mathbb{R})}^{\frac{1}{1+\theta_{1}}}\|g\|_{L^{\frac{1}{\theta_1}}(\mathbb{R})}^{\frac{1}{1+\theta_{1}}}, &  |E_\lambda|\geq2^{j_0}; \\
   |E_\lambda|^{\frac{1-\theta_{1}}{1+\theta_{1}}}\|f\|_{L^{1}(\mathbb{R})}^{\frac{1}{1+\theta_{1}}}\|g\|_{L^{\frac{1}{\theta_1}}(\mathbb{R})}^{\frac{1}{1+\theta_{1}}}, & |E_\lambda|<2^{j_0}.
 \end{cases}
\end{align}

Our next goal is to consider $\Delta_{3}$, a calculation gives
\begin{align*}
\min\left\{2^j,\frac{2^{j(1-\theta_{1})}|E_\lambda|^{\theta_1}}{|\gamma'(2^{j-1})|^{\theta_{1}}},2^{j\theta_1} |E_\lambda|^{1-\theta_1},\frac{|E_\lambda|}{|\gamma'(2^{j-1})|^{\theta_1}}\right\}=
 \begin{cases}
   2^j, &  |E_\lambda|\geq2^j|\gamma'(2^{j-1})|; \\
   \frac{2^{j(1-\theta_{1})}|E_\lambda|^{\theta_1}}{|\gamma'(2^{j-1})|^{\theta_{1}}},&2^j\leq|E_\lambda|<2^j|\gamma'(2^{j-1})|;\\
  \frac{|E_\lambda|}{|\gamma'(2^{j-1})|^{\theta_1}} , & |E_\lambda|<2^j.
 \end{cases}
 \end{align*}
We can then write $\Delta_{3}=\Delta_{3,1}+\Delta_{3,2}+\Delta_{3,3}$, where
$$\begin{cases}
   \Delta_{3,1}:= \sum\limits_{  \substack{j>j_{0}+2   \\|E_\lambda|\geq2^j|\gamma'(2^{j-1})|  }}    2^{{\frac{\alpha}{1+\theta_{1}}}j}
\|f\|_{L^{1}(\mathbb{R})}^{\frac{1}{1+\theta_{1}}}\|g\|_{L^{\frac{1}{\theta_1}}(\mathbb{R})}^{\frac{1}{1+\theta_{1}}}; \\
   \Delta_{3,2}:=\sum\limits_{  \substack{j>j_{0}+2   \\2^j\leq|E_\lambda|<2^j|\gamma'(2^{j-1})|  }}   {\frac{2^{{\frac{\alpha-\theta_{1}}{1+\theta_{1}}}j}|E_\lambda|^{\frac{\theta_{1}}{1+\theta_{1}}}}{|\gamma'(2^{j-1})|^{\frac{\theta_{1}}{1+\theta_{1}}}}} \|f\|_{L^{1}(\mathbb{R})}^{\frac{1}{1+\theta_{1}}}\|g\|_{L^{\frac{1}{\theta_1}}(\mathbb{R})}^{\frac{1}{1+\theta_{1}}};\\
 \Delta_{3,3}:=\sum\limits_{  \substack{j>j_{0}+2   \\|E_\lambda|<2^j}}   {\frac{2^{{\frac{\alpha-1}{1+\theta_{1}}}j}|E_\lambda|^{\frac{1}{1+\theta_{1}}}}{|\gamma'(2^{j-1})|^{\frac{\theta_{1}}{1+\theta_{1}}}}}
 \|f\|_{L^{1}(\mathbb{R})}^{\frac{1}{1+\theta_{1}}}\|g\|_{L^{\frac{1}{\theta_1}}(\mathbb{R})}^{\frac{1}{1+\theta_{1}}}.
 \end{cases}$$
Corresponding to $\Delta_{3,1}$, for any $\varepsilon>0$, it follows from Lemma \ref{lemma 2.1} that
$$\sum\limits_{  \substack{j>j_{0}+2   \\|E_\lambda|\geq2^j|\gamma'(2^{j-1})|  }}    2^{{\frac{\alpha}{1+\theta_{1}}}j}\lesssim
\begin{cases}
   |E_\lambda|^{\frac{\alpha}{(\omega_{\infty,1}-\varepsilon)(1+\theta_{1})}}, & |E_\lambda|\geq2^{j_0+2}|\gamma'(2^{j_0+1})|; \\
  0, & 2^{j_0+2}\leq|E_{\lambda}|<2^{j_0+2}|\gamma'(2^{j_{0}+1})|;\\
  0, &|E_{\lambda}|<2^{j_0+2}.
\end{cases}$$
On the other hand, let us set $\varepsilon\in(0,\omega_{\infty,1})$, and we take that $\alpha<\theta_{1}(\omega_{\infty,1}-\varepsilon)$, then by Lemma \ref{lemma 2.1} one has
$$\sum\limits_{  \substack{j>j_{0}+2   \\2^j\leq|E_\lambda|<2^j|\gamma'(2^{j-1})|  }}   {\frac{2^{{\frac{\alpha-\theta_{1}}{1+\theta_{1}}}j}|E_\lambda|^{\frac{\theta_{1}}{1+\theta_{1}}}}{|\gamma'(2^{j-1})|^{\frac{\theta_{1}}{1+\theta_{1}}}}} \lesssim
\begin{cases}
  |E_\lambda|^{{\frac{\alpha+\theta_1(\omega_{\infty,2}-\omega_{\infty,1}+2\varepsilon)}{(\omega_{\infty,2}+\varepsilon)(1+\theta_{1})}}}, & |E_\lambda|\geq2^{j_0+2}|\gamma'(2^{j_{0}+1})|;\\
  |E_\lambda|^{{\frac{\alpha+\theta_1(\omega_{\infty,2}-\omega_{\infty,1}+2\varepsilon)}{(\omega_{\infty,2}+\varepsilon)(1+\theta_{1})}}}, & 2^{j_0+2}\leq|E_{\lambda}|<2^{j_0+2}|\gamma'(2^{j_{0}+1})|;\\
 0, &|E_{\lambda}|<2^{j_0+2}.
\end{cases}$$
This is the estimate about $\Delta_{3,2}$. Furthermore, for estimating $\Delta_{3,3}$, we need to take $\alpha-1+\theta_{1}>0$. It follows from $\omega_{\infty,1}>1$ that $\alpha-1+\theta_{1}-\theta_{1}(\omega_{\infty,1}-\varepsilon)<0$ for all $\varepsilon\in(0,\omega_{\infty,1}-1]$. From Lemma \ref{lemma 2.1}, we then have
$$\sum\limits_{  \substack{j>j_{0}+2   \\|E_\lambda|<2^j}}  {\frac{2^{{\frac{\alpha-1}{1+\theta_{1}}}j}|E_\lambda|^{\frac{1}{1+\theta_{1}}}}{|\gamma'(2^{j-1})|^{\frac{\theta_{1}}{1+\theta_{1}}}}}\lesssim
\begin{cases}
  |E_\lambda|^{{\frac{\alpha-\theta_1(\omega_{\infty,1}-1-\varepsilon)}{1+\theta_{1}}}},
  & |E_\lambda|\geq2^{j_0+2}|\gamma'(2^{j_{0}+1})|; \\
  |E_\lambda|^{{\frac{\alpha-\theta_1(\omega_{\infty,1}-1-\varepsilon)}{1+\theta_{1}}}},  & 2^{j_0+2}\leq|E_{\lambda}|<2^{j_0+2}|\gamma'(2^{j_{0}+1})|;\\
  |E_\lambda|^{{\frac{1}{1+\theta_1}}}, &|E_{\lambda}|<2^{j_0+2}.
\end{cases}$$
Notice that  $$ |E_\lambda|^{{\frac{\alpha-\theta_1(\omega_{\infty,1}-1-\varepsilon)}{1+\theta_{1}}}}\leq|E_\lambda|^{\frac{\alpha}{(\omega_{\infty,1}-\varepsilon)(1+\theta_{1})}} \leq  |E_\lambda|^{{\frac{\alpha+\theta_1(\omega_{\infty,2}-\omega_{\infty,1}+2\varepsilon)}{(\omega_{\infty,2}+\varepsilon)(1+\theta_{1})}}}$$
holds for $|E_\lambda|\geq2^{j_0+2}$ if $\varepsilon\in(0,\min\{\omega_{\infty,1}, \omega_{\infty,1}-1\})$. Consequently,
\begin{align}\label{eq:4.11}
\Delta_{3}\lesssim
\begin{cases}
   |E_\lambda|^{{\frac{\alpha+\theta_1(\omega_{\infty,2}-\omega_{\infty,1}+2\varepsilon)}{(\omega_{\infty,2}+\varepsilon)(1+\theta_{1})}}}
 \|f\|_{L^{1}(\mathbb{R})}^{\frac{1}{1+\theta_{1}}}\|g\|_{L^{\frac{1}{\theta_1}}(\mathbb{R})}^{\frac{1}{1+\theta_{1}}}, & |E_\lambda|\geq2^{j_0+2};\\
  |E_\lambda|^{{\frac{1}{1+\theta_{1}}}}
\|f\|_{L^{1}(\mathbb{R})}^{\frac{1}{1+\theta_{1}}}\|g\|_{L^{\frac{1}{\theta_1}}(\mathbb{R})}^{\frac{1}{1+\theta_{1}}}, &|E_{\lambda}|<2^{j_0+2}.
\end{cases}
\end{align}

It follows from $\alpha-1+\theta_{1}>0$ and $\theta_1\in (0,1)$ that
$ |E_\lambda|^{{\frac{1}{1+\theta_{1}}}}\leq |E_\lambda|^{\frac{\alpha}{1+\theta_{1}}}\leq |E_\lambda|^{\frac{1-\theta_{1}}{1+\theta_{1}}}$
holds for $|E_{\lambda}|<2^{j_0+2}$. This, combined with \eqref{eq:4.8}-\eqref{eq:4.11}, we conclude that
\begin{align}\label{eq:4.12}
\lambda^{\frac{1}{1+\theta_{1}}}|E_{\lambda}|\lesssim
\begin{cases}
   |E_\lambda|^{\frac{\alpha+\theta_1(\omega_{\infty,2}-\omega_{\infty,1}+2\varepsilon)}{(\omega_{\infty,2}+\varepsilon)(1+\theta_{1})}}
 \|f\|_{L^{1}(\mathbb{R})}^{\frac{1}{1+\theta_{1}}}\|g\|_{L^{\frac{1}{\theta_1}}(\mathbb{R})}^{\frac{1}{1+\theta_{1}}}, & |E_\lambda|\geq2^{j_0+2}; \\
 |E_\lambda|^{\frac{1-\theta_{1}}{1+\theta_{1}}}
\|f\|_{L^{1}(\mathbb{R})}^{\frac{1}{1+\theta_{1}}}\|g\|_{L^{\frac{1}{\theta_1}}(\mathbb{R})}^{\frac{1}{1+\theta_{1}}}, &|E_{\lambda}|<2^{j_0+2}.
\end{cases}
\end{align}
By a calculation, for $\varepsilon$ small enough, we have that $\frac{\omega_{\infty,2}-\omega_{\infty,1}}{\omega_{\infty,2}-1}<\frac{\alpha}{1-\alpha}<\frac{\omega_{\infty,2}}{\omega_{\infty,2}-\omega_{\infty,1}+1}$ implies
\begin{align}\label{eq:4.13}
\max\left\{1-\alpha, \alpha, \frac{\alpha}{\omega_{\infty,1}-\varepsilon}\right\}<\frac{\omega_{\infty,2}-\alpha+\varepsilon}{2\omega_{\infty,2}-\omega_{\infty,1}+3\varepsilon}.
\end{align}
Then, we can take
\begin{align}\label{eq:4.14}
\theta_1\leq \frac{\omega_{\infty,2}-\alpha+\varepsilon}{2\omega_{\infty,2}-\omega_{\infty,1}+3\varepsilon}
\end{align}
with $\varepsilon$ small enough, which further implies that
$$\frac{\alpha+\theta_1(\omega_{\infty,2}-\omega_{\infty,1}+2\varepsilon)}{(\omega_{\infty,2}+\varepsilon)(1+\theta_{1})}\leq \frac{1-\theta_{1}}{1+\theta_{1}},$$
This, together with \eqref{eq:4.12}, gives
\begin{align}\label{eq:4.15}
\lambda|E_{\lambda}|^{\frac{\omega_{\infty,2}+ \theta_1(\omega_{\infty,1}-\varepsilon)-\alpha+\varepsilon}{\omega_{\infty,2}+\varepsilon}}\lesssim \|f\|_{L^{1}(\mathbb{R})}\|g\|_{L^{\frac{1}{\theta_1}}(\mathbb{R})}
\end{align}
for all $0<|E_\lambda|<\infty$, where $\varepsilon$ small enough and $\theta_1$ satisfies \eqref{eq:4.14}. Therefore, we have established the $L^{1}(\mathbb{R})\times L^{\frac{1}{\theta_{1}}}(\mathbb{R})\rightarrow L^{\frac{\omega_{\infty,2}+\varepsilon}{\omega_{\infty,2}+ \theta_1(\omega_{\infty,1}-\varepsilon)-\alpha+\varepsilon}, \infty}(\mathbb{R})$ boundedness for $H_{\alpha,\gamma}$.

It is easy to see that $\frac{\omega_{\infty,2}+ \theta_1(\omega_{\infty,1}-\varepsilon)-\alpha+\varepsilon}{\omega_{\infty,2}+\varepsilon}$ takes the maximal value $\frac{2(\omega_{\infty,2}-\alpha+\varepsilon)}{2\omega_{\infty,2}-\omega_{\infty,1}+3\varepsilon}$ when $\theta_1= \frac{\omega_{\infty,2}-\alpha+\varepsilon}{2\omega_{\infty,2}-\omega_{\infty,1}+3\varepsilon}$. On the other hand, from \eqref{eq:4.13}, it follows $\frac{\alpha}{\omega_{\infty,1}-\varepsilon}<\frac{\omega_{\infty,2}-\alpha+\varepsilon}{2\omega_{\infty,2}-\omega_{\infty,1}+3\varepsilon}$, which equals to $2\alpha+\varepsilon<\omega_{\infty,1}$, we then have $2\alpha<\omega_{\infty,1}$. This, together with $\alpha<1<\omega_{\infty,1}\leq\omega_{\infty,2}$, shows $3\alpha<\omega_{\infty,1}+\omega_{\infty,2}$. Therefore,
$$\frac{\omega_{\infty,2}-\alpha+\varepsilon}{2\omega_{\infty,2}-\omega_{\infty,1}+3\varepsilon}< \frac{\omega_{\infty,2}-\alpha}{2\omega_{\infty,2}-\omega_{\infty,1}}$$
holds for all $\varepsilon$ small enough. Putting things together, we conclude that $H_{\alpha,\gamma}$ maps $L^{1}(\mathbb{R})\times L^{2r( \varepsilon )   }(\mathbb{R})$ to $L^{r(\varepsilon ), \infty}(\mathbb{R})$, where $\frac{1}{r( \varepsilon ) }:=\frac{2(\omega_{\infty,2}-\alpha)}{2\omega_{\infty,2}-\omega_{\infty,1}}-\varepsilon $ with $\varepsilon$ small enough. Therefore, we finish the proof of Proposition \ref{proposition 4.1}.
\end{proof}

\end{document}
